\section{Integrated refinement and numerical evidence}
\label{sec:integrated-requests}

The common constructor preserves the ordinary coefficient engine when the
input already has polynomial logarithmic coefficients. In particular, a
request for Newton inversion of $x+x^2\log x$ retains that method. A
reciprocal or generalized logarithmic coefficient instead selects the
larger algebra of Section~\ref{sec:logarithmic-scales}. Explicit sector and
special-function constructors retain their own cutoff semantics.

\subsection{Three meanings of refinement}

A numeric second argument to \wl{SeriesRefine} is an exact exclusive
cutoff for exponent-truncated objects, or an integer marker depth when
replaying a depth-truncated source. Compatible ordinary inverse objects reuse the states described in
Section~\ref{sec:refinement-state}; logarithmic objects replay their
retained original function at the requested logarithmic or source-power
cutoff. A derived result replays its operation recipe and preserves the
precision limits of its operands.

\begin{lstlisting}
s = AsymptoticInverse[x + x^2, {x, 0}, {y, 3}];
SeriesRefine[s, 6]
SeriesRefine[s, <|"AdditionalBlocks" -> 3|>]
SeriesRefine[s, <|"Target" -> 1/100, "TargetError" -> 10^-30,
  "Interval" -> {1/200, 1/50}|>]
\end{lstlisting}

The additional-block request counts complete nonzero exponent blocks of
an ordinary inverse, including all resonant contributions at each weight.
It advances the existing cutoff by the smallest normalized positive gap
and delegates to the retained-state refinement. A finite exact inverse
returns immediately when no further blocks exist, even if the numerical
block goal is larger. The result records both \wl{"GoalReached"} and
\wl{"ExactTermination"}; these statements are different. A declared
input error or a resource limit can end refinement earlier, with the best
available expansion retained in the failure data.

The tolerance request returns a numerical certificate association, with
\wl{"ResultType" -> "NumericalCertificate"}. It invokes the exact
rational interval operation of Section~\ref{sec:executable-certificates}, retains
the source expansion, and reports a certified rational source-root center.
The tolerance concerns that root even if the retained symbolic object
represents a power observable. It does
not replace an asymptotic Big-$O$ by a numeric tolerance. Absolute and
relative goals have the same interval and zero-root conventions as
\wl{InverseCertificate}. An unknown request key is rejected rather than
silently ignored.

\subsection{Checking observables against the original equation}

\wl{InverseNumericalCheck} accepts ordinary, logarithmic, Fourier, flat,
exact-core and exponential-core inverses. Special-function adapters use
their stable original-function check. The source and target chart routes
retain their exact transformed equations. For an observable with power
$r\ne1$, the numerical reference is
\[
 H_r(x_*)=\begin{cases}(x_*-x_0)^r,&x_0\text{ finite},\\
 x_*^r,&x_0=\pm\infty.
 \end{cases}
\]
For $r=1$, the displayed inverse and the numerical reference both include
the source translation. The selected source sign recovers a root seed
from the observable, including even powers at negative infinity. The
computed root must lie on the selected real source side.

The returned \wl{"ReferenceRoot"} and \wl{"ReferenceObservable"} have
different meanings, and \wl{"Error"} compares the observable with its
finite expansion. \wl{"ForwardResidual"} is evaluated at the recovered
source seed. The old \wl{"ExactInverse"} field remains a compatibility
alias for the numerical reference root; its name is not evidence of an
exact or certified value. Every such comparison is explicitly numerical
evidence. Unknown forward remainders cannot be evaluated and remain
outside the stored equation's numerical scope.

Exact target substitution precedes decimal evaluation. Thus a target
$10^{100}+1/100$ near the limiting value $10^{100}$ is reduced to its exact
distance before finite-precision arithmetic. Rounding the target first
would erase that distance at ordinary working precisions and select the
wrong branch or seed. Transformed targets follow the same rule.

The integration tests compare observable powers with an independent
quadratic radical, retain the lower Lambert branch, resolve flat-sector
differences at high precision, check the original Fourier and special
function equations, and exercise failed precision and branch contracts.
These checks complement the independently generated coefficient and
composition oracles; they do not replace the remainder theorems.
